\section{Related Work}
    \label{sec:related}

\noindent
\textbf{Sequence modeling in echocardiography video segmentation.} 
Echocardiography video segmentation is significantly affected by factors such as speckle noise, probe motion, and rapid non-rigid cardiac deformations.
The standard temporal attention mechanism \cite{NIPS17_Attention,ICLR21_vit,TMI24_DSA} computes the output $o_t$ by calculating the correlation between the full history of key-value pairs $\{(\boldsymbol{k}_\tau, \boldsymbol{v}_\tau)\}_{\tau=1}^t$ and the current query $q_t$.
Although this mechanism can establish global temporal dependencies, its computational complexity grows quadratically with the sequence length.
\emph{Retrieval} paradigms based on a memory bank~\cite{iccv19_stm,neurips21_stm} have been adopted by methods such as XMem~\cite{eccv22_xmem}, XMem++~\cite{iccv25_xmempp}, Cutie~\cite{cvpr24_cutie}, and SAM 2 \cite{ICLR25_SAM2,CVPR25_DAM}.
Similarly, MemSAM~\cite{deng2024memsam} and SAMed-2~\cite{MICCAI25_SAMed2} maintain temporal consistency by filtering and storing reference frames in a memory bank, leveraging feature matching and retrieval.
However, limited by their discrete storage mechanisms, these methods struggle to fully utilize the continuous historical information of videos.
LRM \cite{katharopoulos_transformers_2020,ICML20_ttt,ICML21_fwp,ICLR25_gdn,ICLR26_ttt_LaCT} offer an alternative feature aggregation strategy: compressing the complete history into a fixed-size hidden state matrix $\mathbf{S}_t$, thereby achieving inference with constant time complexity.
EchoVim~\cite{mm25_echovim}, based on Vision Mamba \cite{CLM24_mamba,ICML24_mamba2,eccv24_videomamba}, introduces this type of architecture to improve computational efficiency.
However, its temporal modeling still relies on bidirectional inference dependent on specific frames and Dynamic Inter-frame Temporal Attention between adjacent frames, which essentially remains sparse modeling based on specific physiological priors.
Works targeting long-sequence video modeling (such as LiVOS~\cite{Liu_2025_CVPR} and GDKVM \cite{ICLR25_gdn,ICCV25_gdkvm}) utilize a linear key-value association mechanism to compress temporal history into a time-dependent hidden state $\mathbf{S}_t$.
However, these methods operate in unconstrained Euclidean space, where scalar gating mechanisms act as indiscriminate weight decay on the singular values of $\mathbf{S}_t$.
Over long sequences, this inevitably collapses the state matrix into a low-rank approximation, destroying fine-grained spatial details critical for valvular motion analysis and reducing the effective memory capacity.
Addressing this rank collapse requires fundamentally abandoning the Euclidean formulation in favor of geometric constraints that preserve the spectral structure of memory states.

\noindent
\textbf{Geometric constraints in recurrent optimization.}
Viewing sequence modeling through an optimization lens, methods like Test-Time Training \cite{ICML20_ttt,ICML25_ttt,arxiv25_ttregression} formulate hidden state updates as gradient descent steps on a surrogate objective.
However, because these updates operate purely in unconstrained Euclidean space, the state remains vulnerable to singular value degradation and rank collapse over long sequences.
In parallel, orthogonal constraints have shown benefits in parameter optimization.
Muon \cite{post24_muon,arxiv25_kimimuon,ICLR26_ttt_LaCT} applies Newton-Schulz iterations \cite{SIAM06_Schur,SIAM06_Newton} to orthogonalize the gradient or momentum rather than the weights themselves, finding the steepest descent direction under spectral norm constraints.
While effective for parameter updates, such constraints address the optimization dynamics rather than the inference-time state evolution.
Orthogonal and unitary constraints on state transitions have long been recognized as mechanisms to prevent rank collapse in traditional RNNs \cite{PUR08_Absil,ICML16_urnn,NIPS16_fullurnn}.
Extending this geometric rigor to the matrix-valued states of modern linear associative memory, however, has been hindered by the prohibitive computational cost of exact Singular Value Decomposition (SVD) for enforcing orthogonality.